\chapter*{阅读约定与文档状态}
\addcontentsline{toc}{chapter}{阅读约定与文档状态}

RPA 是 \emph{Reformulation Proof Assembly}；RPX 是 \emph{Reformulation Proof eXecutable}。
\code{.rpa} 是带版本头的严格 JSON 证书，\code{.rpx} 是同一证明对象的顺序二进制编码。
MCM 保存完整的精确有理系数模型。
调用者另行提供源模型、目标模型和所要求的 contract；证书不能自行降低这个要求。
本手册说明当前安全 Rust kernel 如何核验普通有理系数 MILP 的重构关系，
包括完整可行向量、目标关系、投影、双射、精确点见证及严格改进。

\section*{实现基线与协议状态}
\addcontentsline{toc}{section}{实现基线与协议状态}
本手册于 2026 年 10 月 13 日按实现重写，并于 10 月 14 日同步 V9/V10 冻结与统一 FORMAT 规范，取代 10 月 9 日的 D0 设计手册。
实现基线为 \code{crates/core/src/}：数学对象见 \code{model.rs} 与 \code{proof.rs}，
执行器见 \code{verify.rs}，文件编解码见 \code{codec.rs} 与 \code{rpa.rs}，
公开接口见 \code{api.rs}，命令行见 \code{main.rs}。
静态联合扩展分别由 flow、interval、menu、signed、subflow 模块实现；
point、injectivity、level 模块提供对应的精确 contract 检查。

\begin{longtable}{l p{104mm}}
\toprule 对象 & 当前状态 \\\midrule\endfirsthead
\toprule 对象 & 当前状态 \\\midrule\endhead
MCM & 只接受版本 1；与 RPX 的增量版本独立。\\
RPA / RPX 1--8 & 已实现、已冻结；后续版本沿用既有字段与数学意义。\\
RPA / RPX 9 & 已实现、已冻结的目标水平域；2026-10-14 冻结。\\
RPA / RPX 10 & 已实现、已冻结的严格改进、完整像排除、联合 subflow 与显式半空间域；2026-10-14 冻结。\\
D0 / v0.1 & 历史设计来源，既不是可接受的文本版本，也不是当前文件 ABI。\\
\bottomrule
\end{longtable}

实现完成与协议冻结分别登记。项目于 2026-10-14 明确冻结 V9/V10，
沿用既有字段、字节与数学接受条件；不兼容变更须使用新版本。
原 D0 的 section directory、表达式池、隐式 SSA 编号、逐项 mapping/check、
yield/merge 和可执行 theorem plugin 均不属于当前接口。
附录\ref{app:revisions}保留其设计决定与当前实现的对应关系。

\section*{规范性用语与证据}
\addcontentsline{toc}{section}{规范性用语与证据}
“必须”“不得”“拒绝”描述本版 kernel 的格式约束、数学接受条件和公开行为；
“建议”描述证书作者或集成方的工程选择；“可以”描述规则允许的情况。
数学语义、字段顺序、引用可见性和完整义务是规范性内容。
数据结构示意、性能解释、有限测试、历史测量和审阅记录是说明性证据。
一个文件通过汇编只说明能够编码，不等于证明通过验证。

\code{docs/FORMAT.md} 是 MCM 1 与 RPA/RPX 1--10 的唯一完整格式规范；
\code{audit/PROTOCOL_V9_V10_FREEZE.md} 登记冻结决定与验证身份。
各联合扩展的独立数学论证及审阅来源在相关章节和附录列出。
交付 manifest 固定的旧规范原字节保存在 \code{audit/specification-history/}；
历史验收记录保持交付时点的草案状态，不倒填冻结结论。
若发现手册与源代码不同，应将其作为需要修复的文档或实现问题，
不能让证书通过含糊解释、未定义行为或自动放宽检查取得接受结果。

\section*{可信边界}
\addcontentsline{toc}{section}{可信边界}
验证依赖实际完整模型、精确算术、静态推导规则和受审阅的联合存在引理。
可信实现包含 MPS/LP 导入器、MCM/RPX decoder、有理数库、表达式代数、
执行器及构造扩展模型的 builder。builder 的旧 scope 前提检查和模型合法性检查
并不能独立证明其生成的所有新坐标同时可行；该同时存在性质来自各自的可信引理。

RPA 作者、Python 生成器、求解器输出、审计 JSON、文件名和摘要都不能授予数学事实。
哈希用于记录输入来源；kernel 的接受条件仍然检查实际传入模型上的每一个前提和义务。
库的纯验证路径不读取文件系统、环境或时钟；文件操作与计时由 API 外层承担。
测试与独立有限 oracle 是实现证据，不能替代对 Rust 实现及依赖的形式化证明。

\section*{阅读顺序与符合性范围}
\addcontentsline{toc}{section}{阅读顺序与符合性范围}
数学审阅从第\ref{ch:model}章和第\ref{ch:execution}章开始：先确定所要求的命题，
再检查每个完整输出向量及分支、扩展、反证的语义。
第\ref{ch:instructions}章逐项列出推导与终结点的前置条件及拒绝条件；
第\ref{ch:format}章给出全部二进制记录，
第\ref{ch:loader}章说明真实输入方言、资源上限、API 和诊断，
第\ref{ch:assembly}章定义文本语法。
第\ref{ch:examples}章提供可重复运行的完整实例，
第\ref{ch:performance}章区分逻辑工作量、生成成本和进程测量。

\begin{description}
\item[Verifier] 接受所支持版本时，检查完整文件、独立 request、全部前提与实际模型义务；
不得把跳过的分支、失败的规则或资源耗尽报告为 \code{VERIFIED}。
\item[Assembler / disassembler] 负责严格文本与二进制的转换；
不承担独立数学认证，重新编码可以选择最低所需版本。
\item[证书作者] 安排已完成的推导与完整 mapping，保证每个可达叶使用同一个联合输出向量；
不能借用另一方向、兄弟分支或另一文件的局部事实。
\item[集成方] 明确传入完整模型、contract 和资源策略，依据 status 判断结果，
依据具体 contract 解释其数学结论。
\end{description}

本 kernel 不提供 proof search、一般 MILP 求解器、任意跳转、动态 theorem registry、
通用共享表达式 DAG、专用 batch 编码或数值解重建命令。
当前语料的完整验收不意味着任意重构都有紧凑证书；
\code{INVALID} 或 \code{RESOURCE_LIMIT} 也不意味着相应重构关系在数学上不成立。
